\section{Introduction}\label{sec:intro}

The CD index \citep{funk2017dynamic} compares the later papers that cite a focal paper with the later papers that cite
the focal paper's references. Citing papers that ignore the references push the score towards $+1$, and the paper is
called disruptive; citing papers that also cite them push it towards $-1$, and the paper is called consolidating;
papers that cite only the references pull the score towards zero. The index became widely known through studies of
team size \citep{wu2019large} and of a decline in disruptiveness across science and technology \citep{park2023papers}.
An inventory lists 105 studies that bear on whether innovation is becoming less disruptive
\citep{wu2026inventory}, and recommendations have been made for using the index in research evaluation
\citep{chen2025thresholds}. The index is also debated: records without references can drive measured trends
\citep{holst2024artefacts,holst2026partially,park2026reply}, citation inflation biases it
\citep{petersen2024inflation,petersen2025reanalysis}, and its values depend on the data source
\citep{leibel2023overview,xu2025alternative}.

A question that matters for both uses has not been answered: how does the index score papers that are later
retracted? Retractions for paper-mill production have grown quickly: the Retraction Watch snapshot used here records
11,716 such papers, with at most four retraction notices per year before 2020 and 6,782 in 2023 alone. If fabricated
or mass-produced papers scored as disruptive, an evaluation that rewards disruption would also reward them; if they
scored as consolidating, the index might carry a signal that integrity screening could use. The closest prior work
studied other units: questionable publications were less disruptive than comparable publications
\citep{you2022questionable}, and a disruption measure computed for MeSH concepts differed between the concepts used to
annotate retracted and non-retracted papers in most years \citep{copara2024mesh}. Neither compared the scores of the
retracted papers themselves with those of comparable papers, measured before the retraction.

Opposite expectations are plausible. Fabricated or falsified results can report findings that later work cites as new
evidence without engaging the literature the paper built on; such papers would score as disruptive. Paper-mill products
follow templates and standard methods; if they are cited mainly by similar products that cite the same standard works,
they would score as consolidating. The retraction itself may also change how a paper is cited: citations fall sharply
after a retraction \citep{azoulay2015retractions}, retraction notices appear among the citing items
\citep{hsiao2021continued}, and scores computed over windows that include these years may describe the retraction
rather than the paper.

This study answers the question on the open biomedical citation graph. The comparison has three features. First, it is
temporal: retracted papers enter the main analysis only if the retraction came after their five-year scoring window
closed, so the score reflects how the paper was cited while it was still part of the record. Second, each retracted
paper is compared with never-retracted papers from the same journal and publication year. Third, the design, the
reason classes and every pre-planned analysis were written down before any score was computed. Analyses added later
are labelled exploratory; they include interval bounds, a review positive control, a variant of the index without
its $n_k$ term, reference-length diagnostics and a comparison of citing papers around the retraction notice.

Contributions:
\begin{itemize}
\item A paper-level measurement of the CD index for retracted papers scored before their retraction, overall and by
retraction reason (paper mill, misconduct, honest error), for 4,587 papers, of which 4,586 are compared with 31,662
scored controls in 2,239 journal--year strata, with a post hoc extension to 7,349 scored papers, of which 7,348 are compared.
\item Bounds on the small mean difference and no unadjusted top-tier excess in the pre-specified five-year analysis,
alongside a positive-score shift; the post hoc three-year window leans toward a modest top-tier excess.
\item Evidence that the index hides a difference in how retracted papers are cited: papers that cite a retracted paper
are less likely to also cite its references, and at equal reference length retracted papers reach the top tier more often.
\item Evidence that scores including the years after the notice are shifted by notice records and by a gap that widens
around the notice, with a practical recommendation for evaluators.
\item Open code, run records and specification files that let the analysis be repeated on the public data.
\end{itemize}
